% TREŚĆ ROZDZIAŁU: zastąp wskazówki opisem wykonania rozwiązania.
% Nie wklejaj długiego kodu bezpośrednio do pliku .tex. Umieść go w katalogu
% code/ rozdziału i wczytaj przez \lstinputlisting — przykład jest na końcu.
\chapter{Realizacja projektu}
\label{chap:realizacja}

\section{Środowisko i narzędzia}
Wymień tylko te wersje narzędzi, bibliotek, urządzeń i konfiguracji, które są istotne dla odtworzenia projektu.

\section{Implementacja lub wykonanie}
Opisz najważniejsze etapy budowy rozwiązania i decyzje podjęte podczas realizacji. Oddziel elementy wykonane samodzielnie od gotowych komponentów.

\section{Instrukcja uruchomienia}
Podaj minimalny, jednoznaczny zestaw kroków potrzebnych do uruchomienia, zbudowania albo użycia rozwiązania.

% Kod przechowuj w katalogu code/ odpowiedniego rozdziału i wczytuj np.:
% \lstinputlisting[style=project-python,caption={Opis kodu},label={lst:przyklad}]{chapters/04-realizacja/code/program.py}
